% Programming with InitechBase, Introduction (PUB-REF-0001 AN-21 F-ATE:
% what the section is, who should read it, chapter-by-chapter list,
% Conventions Used in This Manual; annex Sec 1, 8.1).  Folios PI-n.
\SectionIntroduction{Introduction}
\Gates{test-pubs-build}
\doctitle{Programming with InitechBase} teaches you to write programs: files of
InitechBase commands that do a task from beginning to end whenever you run
them.  It is written for readers who have worked through
\doctitle{Learning InitechBase} and know the commands at the dot prompt.  You
need no experience of programming.

The section builds one application, the Initech Accounting System, a chapter
at a time.  The system keeps the accounts receivable ledger of Initech
Systems Corporation in the database \file{INVOICE.DBF}: one record for each
invoice, with its number, the customer, the amount, the due date and whether
it has been paid.  Its programs list and find invoices, enter new ones, print
the invoice register, age the open invoices at the end of the month, and post
the monthly finance charges.  Every program is on the Sample Programs and
Utilities diskette, and every program is printed in full in this section.

For the exact form of a command or function, this section refers you to
\doctitle{Using InitechBase} in Volume One.

\section{How This Section Is Organized}
\SideHead{Chapter 1}Programming in a Nutshell: what a program is, and the first
program of the system.

\SideHead{Chapter 2}Program Structure and Flow: sequence, decisions and loops.

\SideHead{Chapter 3}Memory Variables: totals, types and the life of a memory
variable.

\SideHead{Chapter 4}Procedures and Parameters: dividing a program into parts.

\SideHead{Chapter 5}Functions for Fields and Memory Variables: building and
formatting values.

\SideHead{Chapter 6}Communicating with the User: questions and answers.

\SideHead{Chapter 7}Working with the Database: finding, adding and changing
records in a program.

\SideHead{Chapter 8}Printing: the invoice register.

\SideHead{Chapter 9}The Aging Report: the month-end aging of the open invoices.

\SideHead{Chapter 10}Finance-Charge Posting: the monthly finance charges and
the rounding adjustment.

\SideHead{Chapter 11}The Main Menu: the program the clerk uses every day.

\SideHead{Chapter 12}Handling Errors: what a program does when a command
fails.

\SideHead{Chapter 13}Testing and Debugging: finding the mistake in a program.

\SideHead{Chapter 14}Working with Two Databases: a lookup in a second work
area.

\SideHead{Chapter 15}Reports from Programs: pages, headings and page numbers.

\SideHead{Chapter 16}Working with Dates: years of service.

\SideHead{Chapter 17}Working with Characters: names for letters.

\SideHead{Chapter 18}Changing Many Records: raises by department.

\SideHead{Chapter 19}Memo Fields in Programs: the remarks of the roster.

\SideHead{Chapter 20}Order and Filters in Programs: two orders and a filter.

\SideHead{Chapter 21}Searching: a search loop with \cmd{LOCATE}.

\SideHead{Chapter 22}Putting It Together: the Initech Accounting System as a
whole, and the month-end procedure.

\section{Conventions Used in This Manual}
\SideHead{Programs}Program listings are printed in bold on a shaded panel,
line for line as they are stored on the sample diskette.

\SideHead{Examples}The panels of examples show the screen: the dot prompt, a
period and a blank, the lines you type in bold, and what InitechBase displays.
Every example was run against a fresh copy of the sample diskette; your
results agree with the book if your files are as they were on the diskette.

\SideHead{Names}The programs of this section begin the names of their memory
variables with M, so that a memory variable is never taken for a field.

\SideHead{Dates}The examples were run on the 1st of March 1992, the date
\fn{DATE()} returns in them.  The accounting programs take their reporting
date from a memory variable instead, so that a month can be closed again
later.
